\documentclass[11pt]{article}
\usepackage[margin=1in]{geometry}
\usepackage{amsmath,amssymb,amsthm,hyperref}
\newtheorem{proposition}{Proposition}
\newtheorem{lemma}{Lemma}
\newcommand{\C}{\mathbb C}
\newcommand{\R}{\mathbb R}
\newcommand{\Pcal}{\mathcal P}
\newcommand{\norm}[1]{\left\lVert#1\right\rVert}
\newcommand{\abs}[1]{\left|#1\right|}
\title{Joint endpoint estimates and localized two-window projection transport}
\author{Split-Zero research continuation}
\date{19 September 2026}
\begin{document}
\maketitle

\paragraph{Source and precise scope.}
This note verifies and supplies finite constants for the deductions labelled
J1--J10 and S15--S16 in the supplied \texttt{COMPLETE\_PROOFS.md}, dated
18 September 2026.  The source bytes have SHA-256
\texttt{e4e54261c91ae899a78956383482650f88933024f5f3cec6b41332ceeb5a4172}.
The endpoint block is at lines 799--966; the two-window block is at
lines 746--770.  This note retains every singular coefficient, the original
Lebesgue product measure, the full polynomial relation, and both boundary
windows.  The incoming endpoint decomposition, the actual tilt envelope,
and the old window coefficients are explicitly identified as input results.
The named earlier files \texttt{TWO\_WINDOW\_SAME\_CLASS\_DEFECT.md} and the
theta-endpoint note were not located in the bounded intake.  Their numerical
window coefficient and full theta existence theorem are therefore not
claimed as independently recovered here.  The finite deductions below are
proved from the actual formulas recorded in the supplied source; no missing
allocation or projected arithmetic sign is assumed.

\section{The entire remainder has an explicit uniform norm allowance}

Write the original quartet as
\[
\mathcal Z=\{1/2+\epsilon\delta+i\epsilon'\gamma:
                 \epsilon,\epsilon'\in\{-1,1\}\},
\qquad 0<\delta<1/2,\quad\gamma>2,
\]
with common multiplicity \(m\), \(d=m-1\), and retain the given inverse
primitive
\[
\psi(x)=s_h(x)+R_h(x),\qquad
s_h(x)=\sum_{\rho\in\mathcal Z}\sum_{j=0}^{d}
a_{\rho,j}x^{-\rho}\frac{\log(1/x)^j}{j!}.
\tag{JL1}
\]
Here \(R_h\) is the actual entire even remainder, not a freely chosen
regularization.  If
\(R_h(x)=\sum_{n\ge0}r_{2n}x^{2n}\), put
\[
B_h=\sum_{n\ge0}|r_{2n}|<\infty,\qquad
C_h^{\rm reg}=B_h^2+
2B_h\sum_{\rho,j}
\frac{|a_{\rho,j}|}{(1-\Re\rho)^{j+1}}.
\tag{JL2}
\]
The sum defining \(B_h\) converges absolutely because the power series
is entire.  One can alternatively use the possibly smaller
\(\sup_{0\le x\le1}|R_h(x)|\).  These constants refer to that same remainder
and leave all singular coefficients unchanged.

\begin{proposition}
For every \(L\ge0\),
\[
\left|\int_{e^{-L}}^1|\psi(x)|^2\,dx
      -\int_{e^{-L}}^1|s_h(x)|^2\,dx\right|
\le C_h^{\rm reg}.
\tag{JL3}
\]
\end{proposition}
\begin{proof}
The integrand difference is
\(2\Re(s_h\overline{R_h})+|R_h|^2\).  Since
\(\Re\rho<1\),
\[
\int_0^1 x^{-\Re\rho}\frac{\log(1/x)^j}{j!}\,dx
=\frac1{(1-\Re\rho)^{j+1}}.
\]
This identity follows from \(x=e^{-t}\) and \(j\) integrations by parts
in \(\int_0^\infty e^{-(1-\Re\rho)t}t^j\,dt\).
Take absolute values of each of the finitely many singular terms,
use \(|R_h|\le B_h\), and enlarge the integration interval to \((0,1)\).
This proves (JL3).  In particular the cross terms with the entire
remainder do not grow with the endpoint cutoff.
\end{proof}

\section{All factorials in the oscillating endpoint leading term}

Let
\(\rho_+=1/2+\delta+i\gamma\) and
\(\mathcal R=\{\rho_+,\overline{\rho_+}\}\).
The local endpoint coefficient convention in the supplied source is
\[
a_{\rho,j}=[z^{d-j}]
\frac{p_\delta(\rho+z)}{2z_\rho(z)},\qquad
p_\delta(s)=(s-1/2)^2+\gamma^2-\delta^2,\qquad
\zeta(\rho+z)=z^m z_\rho(z).
\tag{JL4}
\]
Thus
\[
a_{\rho_+,d}
=\frac{2i\delta\gamma}{2z_{\rho_+}(0)}
=\frac{m!\,i\delta\gamma}{\zeta^{(m)}(\rho_+)}
=:A_h\ne0,\qquad
a_{\overline{\rho_+},d}=\overline{A_h}.
\tag{JL5}
\]
This is the top Laurent coefficient in the recorded convention; no
phase has been made real by a coordinate change.

Put \(I_h(L)=\int_{e^{-L}}^1|\psi|^2\,dx\).
The change of variables \(x=e^{-t}\) gives the exact finite identity
\[
I_h^{\rm sing}(L)=
\sum_{\rho,\sigma\in\mathcal Z}\sum_{j,l=0}^d
\frac{a_{\rho,j}\overline{a_{\sigma,l}}}{j!\,l!}
E_{j+l}(\rho+\bar\sigma-1,L),\qquad
E_p(z,L)=\int_0^L e^{zt}t^p\,dt.
\tag{JL6}
\]
There is no division by \(z\) in this definition, so the terms with
\(z=0\), including the reflected-root pairs, remain present.
For \(z\ne0\), repeated integration by parts gives
\[
E_p(z,L)=e^{zL}\sum_{r=0}^p
\frac{(-1)^r p!L^{p-r}}{(p-r)!z^{r+1}}
-\frac{(-1)^p p!}{z^{p+1}}.
\tag{JL7}
\]

For every ordered right/right pair and \(p=j+l\) put
\[
\begin{split}
A_1={}&\sum_{\rho,\sigma\in\mathcal R}\sum_{j,l=0}^d
\frac{|a_{\rho,j}a_{\sigma,l}|}{j!\,l!}
\sum_{\substack{0\le r\le j+l\\(j,l,r)\ne(d,d,0)}}
\frac{(j+l)!}{(j+l-r)!\,|\rho+\bar\sigma-1|^{r+1}},\\
A_2={}&C_h^{\rm reg}+
\sum_{\rho,\sigma\in\mathcal R}\sum_{j,l=0}^d
\frac{|a_{\rho,j}a_{\sigma,l}|}{j!\,l!}
\frac{(j+l)!}{|\rho+\bar\sigma-1|^{j+l+1}},\\
A_3={}&
\sum_{(\rho,\sigma)\in\mathcal Z^2\setminus\mathcal R^2}
\sum_{j,l=0}^d
\frac{|a_{\rho,j}a_{\sigma,l}|}{j!\,l!\,(j+l+1)}.
\end{split}
\tag{JL8}
\]
These are exactly J2 with (JL2) as a fully displayed valid regular
allowance.  Define
\[
F_h(L)=\frac1{(d!)^2}
\left(\frac{|A_h|^2}{\delta}
+\Re\frac{A_h^2e^{2i\gamma L}}{\delta+i\gamma}\right),
\qquad
f_*=\frac{|A_h|^2}{(d!)^2}
\left(\frac1\delta-\frac1{\sqrt{\delta^2+\gamma^2}}\right)>0.
\tag{JL9}
\]

\begin{proposition}
For every \(L\ge1\),
\[
F_h(L)\ge f_*,
\qquad
\left|\frac{I_h(L)}{e^{2\delta L}L^{2d}}-F_h(L)\right|
\le \frac{A_1}{L}+(A_2+A_3L)e^{-2\delta L}.
\tag{JL10}
\]
\end{proposition}
\begin{proof}
For the two equal right/right roots, the \(j=l=d,r=0\) terms in
(JL7) sum to
\(|A_h|^2/(\delta(d!)^2)\) after dividing by the displayed scale.
The other two ordered right/right pairs sum to
\(\Re(A_h^2e^{2i\gamma L}/(\delta+i\gamma))/(d!)^2\).
These are precisely the terms in (JL9).
Every remaining upper endpoint term in (JL7), after that division,
contains \(L^{j+l-r-2d}\) with exponent at most \(-1\).
The sum of its absolute coefficients is \(A_1\).
The lower endpoint terms and (JL3) are bounded by
\(A_2e^{-2\delta L}\), because \(L^{-2d}\le1\).
For every pair outside \(\mathcal R^2\), \(\Re z\le0\), and therefore
\[
|E_{j+l}(z,L)|\le \frac{L^{j+l+1}}{j+l+1}.
\]
Its scaled contribution is at most the corresponding coefficient in
\(A_3Le^{-2\delta L}\).
The lower bound for \(F_h\) follows by taking the modulus of its
oscillating summand.  This proves all assertions, retaining all ordered
cross terms.
\end{proof}

\paragraph{Joint tensor/cutoff error.}
Put
\[
\eta_h(L)=
\frac{A_1/L+(A_2+A_3L)e^{-2\delta L}}{f_*}.
\]
If \(\eta_h(L)\le1/2\), the actual ratio
\(I_h(L)/(e^{2\delta L}L^{2d}F_h(L))\) lies in
\([1-\eta_h(L),1+\eta_h(L)]\).
For real \(|x|\le1/2\), integration of \(1/(1+t)\) from \(0\) to \(x\)
proves \(|\log(1+x)|\le2|x|\).  Consequently
\[
\left|\log
\frac{I_h(L)^k}{e^{2\delta kL}L^{2dk}F_h(L)^k}\right|
\le2k\eta_h(L),
\tag{JL11}
\]
with relative error at most \(e^{2k\eta_h(L)}-1\).
At \(L=k^2\), the exponent is \(O_h(k^{-1})\), so
\[
I_h(k^2)^k=e^{2\delta k^3}k^{4dk}F_h(k^2)^k
\bigl(1+O_h(k^{-1})\bigr).
\tag{JL12}
\]
When \(m=1\), the empty sum \(A_1\) is zero.  At
\(L=(\log k)/\delta\ge1\), the exponent in (JL11) is at most
\[
\frac{2}{f_*k}\left(A_2+\frac{A_3}{\delta}\log k\right).
\tag{JL13}
\]
Fubini's theorem gives exactly \(I_h(L)^k\) for the fixed original
product representative on \([e^{-L},1]^k\).
These estimates concern that specified representative.  They do not
evaluate a minimum over added source corrections.

\section{The literal diagonal map and the terminal factorial}

For \(b_\rho x^{-\rho}\), \(\Re\rho=1/2+\delta\), the unchanged product
measure gives
\[
N_{\rm prod}(k,L)=|b_\rho|^{2k}
\left(\frac{e^{2\delta L}-1}{2\delta}\right)^k.
\tag{JL14}
\]
The literal diagonal map substitutes all \(x_i=x\), hence its value is
\(b_\rho^kx^{-k\rho}\).  Its norms in the two expressly specified measures are
\[
\begin{split}
N_{\rm diag,dx}(k,L)&=|b_\rho|^{2k}
\frac{e^{(k-1+2k\delta)L}-1}{k-1+2k\delta},\\
N_{\rm diag,k}(k,L)&=|b_\rho|^{2k}
\frac{e^{2k\delta L}-1}{2k\delta}
\quad\text{in }x^{k-1}\,dx.
\end{split}
\tag{JL15}
\]
Integrating the respective powers of \(x\) proves these formulas.
Dividing (JL14) by the second formula gives the exact discrepancy
\[
\frac{N_{\rm prod}}{N_{\rm diag,k}}
=k(2\delta)^{1-k}
\frac{(1-e^{-2\delta L})^k}{1-e^{-2k\delta L}}.
\tag{JL16}
\]

The common terminal factorial can also be checked without dropping the
original unit.  In one local algebra \(\C[z]/(z^m)\), the actual terminal
vector is \(z^d\); its endpoint coefficient is
\(b_\rho=1/(2z_\rho(0))\) in convention (JL4).
In the extremal ordered tuple, put \(n=kd\) and
\(N_\Sigma=z_1+\cdots+z_k\).  The multinomial theorem and \(z_i^{d+1}=0\)
give the exact polynomial identity
\[
N_\Sigma^{kd}=\frac{(kd)!}{(d!)^k}\prod_{i=1}^k z_i^d.
\tag{JL17}
\]
Multiplying by the full \(U_h^{\otimes k}\) kills only its positive-degree
terms in this terminal vector, so its endpoint is
\[
\frac{(kd)!}{(d!)^k}\,U_h(\rho)^k b_\rho^k
\prod_{i=1}^k x_i^{-\rho}.
\tag{JL18}
\]
Both product and diagonal squared norms therefore acquire
\(\left|(kd)!/(d!)^k\right|^2|U_h(\rho)|^{2k}\).
This verifies the sentence following J8 for that exact terminal vector.
No factorial is absorbed into an unnamed choice of basis.

\section{Uniform propagation of moments to the closed tilt interval}

Retain the actual nonnegative one-packet density \(w_h\) and its stated
envelope
\[
e^{\alpha|y|}w_h(y)\le C_h^{\rm tilt}(1+|y|)^{-p},
\qquad \alpha=\pi/2,\quad p=8m-9/2>3.
\tag{JL19}
\]
Let \(\mu_n=\int_\R y^nw_h(y)\,dy\), without changing the source mass.
For \(j=0,1,2\),
\[
\sum_{r\ge0}\frac{\alpha^r}{r!}
\int_\R |y|^{r+j}w_h(y)\,dy
=\int_\R |y|^je^{\alpha|y|}w_h(y)\,dy<\infty.
\tag{JL20}
\]
Tonelli's theorem applies to this nonnegative majorant.  It proves
uniform absolute convergence on \(0\le\theta\le\alpha\), interchange
with the signed integrands, and continuity of the derivatives at
the endpoint:
\[
M_h^{(j)}(\theta)
=\int_\R y^je^{\theta y}w_h(y)\,dy
=\sum_{r\ge0}\mu_{r+j}\frac{\theta^r}{r!}.
\tag{JL21}
\]

For an integer cutoff \(L\ge0\), let \(n=L+1\) and
\(T_L=n/(2e\alpha)\).  For \(z\ge0\) put
\[
\mathcal Q_n(z)=e^{-z}\sum_{r=n}^{\infty}\frac{z^r}{r!}.
\]
For \(0<z\le n/(2e)<n\), multiplication of each term by
\((n/z)^{r-n}\ge1\) yields
\[
\mathcal Q_n(z)
\le (z/n)^n e^{n-z}
\le(ez/n)^n\le2^{-n}.
\tag{JL22}
\]
For \(z=0\) the sum is zero; globally it is at most one.
The absolute tail of (JL21) is bounded by
\[
\int_\R |y|^jw_h(y)e^{\alpha|y|}
\mathcal Q_n(\alpha|y|)\,dy.
\]
Use \(2^{-n}\) on \(|y|\le T_L\), and one elsewhere.
Since
\[
\int_0^\infty y^j(1+y)^{-p}\,dy
=\operatorname B(j+1,p-j-1),\qquad
\int_{T_L}^\infty y^j(1+y)^{-p}\,dy
\le\frac{(1+T_L)^{j-p+1}}{p-j-1},
\]
the resulting explicit bound is
\[
\sup_{0\le\theta\le\alpha}
\left|M_h^{(j)}(\theta)
-\sum_{r=0}^{L}\mu_{r+j}\frac{\theta^r}{r!}\right|
\le2C_h^{\rm tilt}\left[
2^{-(L+1)}\operatorname B(j+1,p-j-1)
+\frac{(1+T_L)^{j-p+1}}{p-j-1}\right].
\tag{JL23}
\]
The beta identity follows by \(u=y/(1+y)\) in its defining integral.
Thus no fourth tilted moment is used.  These formulas verify J9--J10
but do not by themselves verify the separately cited arithmetic
moment recurrence or evaluate a growing projection-kernel integral.

\section{The localized ambient space really contains the window projections}

Here restore the simple target order in S15: \(a=k+4\),
\(a\equiv1\bmod4\), \(q=(a+1)^2\), \(c=a/2\), and
\[
\chi(S)=\prod_{r,s=0}^{a}
\bigl[S-c-(2r-a)\delta-i(2s-a)\gamma\bigr].
\tag{JW1}
\]
The actual eigenvalue \(\lambda\) is one of these roots.
Set \(D=2q+1\),
\[
\mathcal W=\frac{\chi(S)}{S-\lambda}\Pcal_{D-q+1},
\quad
\mathcal R_N=\chi\Pcal_{N-q}
\quad (\Pcal_j=\{0\}\text{ for }j<0).
\tag{JW2}
\]
All norms are integrals on the original line \(S=c+iy\).
Let \(\omega\) be the actual arithmetic density and
\(\widetilde\omega\) its specified central replacement at order \(a\),
center \(c\), and interval \([-a\mu_h,a\mu_h]\).
The central-localization calculation C3--C13 gives on this \emph{entire}
fixed space the relative form bound
\[
|(\widetilde\omega-\omega)(f,f)|\le\tau\|f\|_\omega^2
\quad(f\in\mathcal W),\qquad 0\le\tau<1.
\tag{JW3}
\]
For clarity, the valid finite parameter is
\[
\tau=B_a\min\{1,\mathfrak e_{a,1,2q+1}(\omega)\},
\qquad
B_a=e^{H_*}-1,
\tag{JW4}
\]
with \(H_*\) and \(\mathfrak e\) the actual constants in C3 and C13.
The guards of those expressions, including \(R<q\) and \(a\mu_h\le q/2\),
are retained.  Self-adjointness and polarization of (JW3) yield
\[
|(\widetilde\omega-\omega)(f,g)|
\le\tau\|f\|_\omega\|g\|_\omega.
\tag{JW5}
\]

\paragraph{Exact fixed relation band.}
Let \(S_Nv\) be the original attained polynomial representative of \(v\)
at cutoff \(N\).  It is orthogonal to \(\mathcal R_N\).
Every representative of the chosen eigenclass belongs to \(\mathcal W\);
all its differences belong to \(\mathcal R_N\).
For \(n=q-1\) or \(n=q\), put
\[
F_n=S_{n+q}v-S_nv,\qquad
\mathcal B_n=\mathcal R_{n+q}\cap\mathcal R_n^{\perp_\omega}.
\tag{JW6}
\]
Then \(F_n\in\mathcal B_n\), and the \(\omega\)-orthogonal projector onto
this band is \(B_n=P_{\mathcal R_{n+q}}-P_{\mathcal R_n}\).
The analogous statement holds in the replaced metric.
The \emph{full} section-difference map from the cyclic module to
\(\mathcal B_n\) has rank \(q\), not merely on the single eigenline.
Indeed
\(\dim\mathcal R_{n+q}-\dim\mathcal R_n=q\).
If the section difference of any class is zero, its common representative
\(P\), of degree at most \(n\), is orthogonal to every high relation.
In particular it is orthogonal to \(\chi P\), whose degree is at most
\(n+q\).

To check the sign used here, pair each root in (JW1) with the root
having opposite real displacement and the same imaginary displacement.
On \(S=c+iy\), each pair contributes
\(-[(y-(2s-a)\gamma)^2+(2r-a)^2\delta^2]\).
There are \(q/2\) pairs, and \(q/2\) is even because \(a\equiv1\bmod4\).
Every \((2r-a)\delta\) is nonzero because \(a\) is odd.
Thus \(\chi(c+iy)>0\) for every real \(y\).  For \(P\ne0\),
\(\int\chi(c+iy)|P(c+iy)|^2d\omega(y)>0\), contradicting orthogonality.
The map is injective and hence fills the band.

\begin{lemma}[Explicit fixed-subspace projection bound]
For any fixed \(\mathcal R\subseteq\mathcal W\), let \(P,\widetilde P\)
be its orthogonal projectors in the two metrics.  Put
\[
d_\tau=\frac{\tau}{1-\tau},\qquad
g_\tau=\sqrt{\frac{1+\tau}{1-\tau}}.
\]
Then, as operators on \((\mathcal W,\omega)\),
\[
\|\widetilde P-P\|\le d_\tau,\qquad
\|\widetilde P\|\le g_\tau.
\tag{JW7}
\]
For the window band,
\[
\|\widetilde B_n-B_n\|\le2d_\tau,\qquad
\|\widetilde B_n\|,\ \|I-\widetilde B_n\|\le g_\tau.
\tag{JW8}
\]
\end{lemma}
\begin{proof}
For \(w\in\mathcal W\), put \(r=w-Pw\) and
\(h=\widetilde Pw-Pw\in\mathcal R\).  The original orthogonality gives
\(\langle r,h\rangle_\omega=0\); the new orthogonality gives
\(\langle r-h,h\rangle_{\widetilde\omega}=0\).
Consequently
\[
(1-\tau)\|h\|_\omega^2
\le \|h\|_{\widetilde\omega}^2
= (\widetilde\omega-\omega)(r,h)
\le\tau\|r\|_\omega\|h\|_\omega
\le\tau\|w\|_\omega\|h\|_\omega.
\]
This proves the difference bound.  For the second bound use
\(\|\widetilde Pw\|_{\widetilde\omega}\le
\|w\|_{\widetilde\omega}\) and (JW3).
Subtracting the two fixed nested relation projectors proves the first
part of (JW8).  Their difference and its complement are orthogonal
projectors in the new metric, proving its second part.
\end{proof}

\paragraph{Sections, multiplication, and both projected errors.}
The same argument in a fixed affine observation fibre proves
\[
\|\widetilde S_Nv-S_Nv\|_\omega
\le d_\tau\|S_Nv\|_\omega.
\tag{JW9}
\]
Put \(E_n=\|S_nv\|_\omega^2\).
Nested minima and orthogonality give
\[
\|S_{n+q}v\|_\omega\le\sqrt{E_n},\qquad
\|F_n\|_\omega^2=E_n-E_{n+q}\le E_n,\qquad
\|\widetilde F_n-F_n\|_\omega\le2d_\tau\sqrt{E_n}.
\tag{JW10}
\]
Multiplication by \(S-\lambda\) takes all these section differences
into \(\chi\Pcal_{D-q}\subseteq\mathcal W\), with degree at most \(D\).
Let \(M_a\) be a valid operator allowance for this multiplication on
the input face in the original metric.  The actual original
Gamma comparison C2 gives the following fully explicit choice:
\[
M_a=
\sqrt{\frac{b_a}{a_{a,\,2q+1}}}
\bigl(4q+2+|\lambda-c|\bigr).
\tag{JW11}
\]
In fact the incoming Gamma multiplication bound on degree \(2q\)
is \(\|yP\|_\sigma\le2(2q+1)\|P\|_\sigma\);
the triangle inequality for
\((S-\lambda)P=(iy-(\lambda-c))P\), the upper comparison for the output,
and the lower comparison for the input give (JW11).
Thus this bound uses only the displayed C2 provider.  It does not
silently replace \(b_a/a_{a,2q+1}\) by a fixed constant.
If the norms are written with the unscaled density \(m_a\) instead of
\(M_h(\alpha)^{-a}m_a\), both source/Gamma comparison constants acquire
the same factor \(M_h(\alpha)^a\).  Their ratio in (JW11) is unchanged;
all section energies in the formulas still carry their original mass.

Set \(w_n=(S-\lambda)F_n\) and
\(\widetilde w_n=(S-\lambda)\widetilde F_n\).  From (JW10)--(JW11),
\[
\|\widetilde w_n-w_n\|_\omega
\le2d_\tau M_a\sqrt{E_n}.
\tag{JW12}
\]
Equations (JW8) and (JW12), with the triangle inequality, now give
\[
\begin{split}
\|\widetilde B_n\widetilde w_n-B_nw_n\|_\omega
&\le2d_\tau(g_\tau+1)M_a\sqrt{E_n},\\
\|(I-\widetilde B_n)\widetilde w_n-(I-B_n)w_n\|_\omega
&\le2d_\tau(g_\tau+1)M_a\sqrt{E_n}.
\end{split}
\tag{JW13}
\]
Every vector and every projection in these inequalities is in
\(\mathcal W\).  No restricted pair-metric estimate is applied to
an arbitrary projector on the full cyclic quotient.

For a squared-norm version put
\(\varepsilon_\tau=2d_\tau(g_\tau+1)\).
For either projected vector \(V_n\) and its corresponding
\(\widetilde V_n\),
\[
\left|\|\widetilde V_n\|_{\widetilde\omega}^2
              -\|V_n\|_\omega^2\right|
\le
\left[\tau(1+\varepsilon_\tau)^2
      +2\varepsilon_\tau+\varepsilon_\tau^2\right]M_a^2E_n.
\tag{JW14}
\]
To see this, use \(\|V_n\|_\omega\le M_a\sqrt{E_n}\),
(JW13), and (JW3), then expand the difference of the two squared
norms.  This includes the cross terms.
Also, by comparing each of the two attained minima,
\[
\left|\|\widetilde F_n\|_{\widetilde\omega}^2
              -\|F_n\|_\omega^2\right|
\le\tau(E_n+E_{n+q})\le2\tau E_n.
\tag{JW15}
\]

\paragraph{Consequence for the inherited leading coefficients.}
The supplied original two-window result records
\(\|F_n\|_\omega^2/E_n\to1\) and the total coefficient
\(16/\pi^2\) after division by \(q^2\|F_n\|_\omega^2\);
it records that this coefficient is tangential at \(n=q-1\)
and normal at \(n=q\).  Equations (JW14)--(JW15) prove stability
of each of those exact existing coefficients under the specified
central replacement.  Indeed the displayed C2 constants give
\[
\log(b_a/a_{a,2q+1})=O_h(a+\log q),\qquad
\log\tau\le-2(q-1)\log(q/a)+O_h(q+a+\log q).
\tag{JW16}
\]
Since \(q=(a+1)^2\) and \(|\lambda-c|\le
a\sqrt{\delta^2+\gamma^2}\), it follows that
\(\tau M_a^2/q^2\to0\).  Hence the right side of (JW14),
divided by \(q^2E_n\), tends to zero for each window.
The denominator perturbation (JW15) also tends to zero relatively.
This proves the transfer without needing an unlocated sharp
\(O_h(q)\) multiplication theorem.

The bare rate \(O_h(\tau q\sqrt{E_n})\) stated in S16 should therefore
be replaced in the received source by the finite (JW13), unless its
separately named sharper multiplication provider is supplied.
Its asserted leading-coefficient stability is already justified by
the weaker explicit rate.  The numerical coefficient of the old
two-window theorem remains an inherited value: this proof does not
claim to have independently calculated that absent source.

\paragraph{Integration boundary.}
The results verified here neither calculate the seven absolute native
allocations nor the individual three-column projected arithmetic signs.
They provide a quantitative fixed-representative endpoint limit,
the exact diagonal norm discrepancy, uniform fixed-source moment input,
and complete transport of two specified boundary-window projections.
All original allocation and same-class response calculations remain in
their original metrics.

\end{document}
